% !TeX root = ../papers/nc-p3.tex
\section{両切断に階数1の色が三つある全窓}
\label{sec:klein-three-lines}
\begin{theorem}\label{thm:klein-three-lines}
正則表現型の四元群窓で隣接する二切断の色階数がともに $(3,1,1,1)$ の置換であるとする。
各切断の階数3の色を位置別の文字平行移動で零に取る。
このとき標準AS窓は、各位置での対角基底変更の後に、次のテンソルとその四元群二コサイクル捻りで尽きる。
\[
 W_{ijkl}={\bf1}_{i+j+k+l=0}
 -\sum_{a\in K}p_a^{-1}{\bf1}_{i=j=k=l=a},
 \qquad \sum_a p_a=1,\quad \prod_a p_a\ne0 .
\]
表示の全点が左右Koszul AS正則で、最小周期は1である。両表示の間に添字保持同型はない。
対称表示の反対切断の階数は6である。捻った表示では
$\sum_a1/p_a\ne4$ のとき16、$(-1,1,1,1)/2$ の置換四点では14、
それ以外の $\sum_a1/p_a=4$ の点では15である。
\end{theorem}
\begin{proof}[零成分を持つ台の検査]
$M_b(i,k)=w_{i,i+b,k,k+b}$ とし、一位置巡回したものを $N_e$ と書く。
階数1の非零行列の台は非空な行集合と列集合の直積であり、225通りある。
$N_1,N_2,N_3$ の三長方形を選び、$M_1,M_2,M_3$ のそれぞれについてこれと
交わる12セルに一致する長方形を全て取る。選択後に未定なのは $w_{iiii}$ の4セルだけである。
$M_0,N_0$ それぞれに大きさ3のマッチングがあり、完全マッチングがちょうど1個でないことを要求すると
39461472台となる。これは階数3の必要条件だけを使う段階である。

零行・列から四層の自由辺グラフを作ると、強連結成分内に分岐を持たない台は5160個である。
このうち16個は六つの長方形が全て四行四列である。残る5144個を同時文字置換、巡回、反転でまとめると
52軌道となる（四文字の任意の置換がアフィン変換である）。
行集合 $I$ への射影階数を、色階数とその行集合上の最大マッチング数の小さい方で抑える。
前節までの連続純関係検査と一文字商の三次元数下界
（\path{sections/klein_single_missing_supports.tex}）をこの上界で実行すると、52軌道は全て排除される。
台のビット番号は $16(i+j)+4i+k$（$i+j$ は四元群の和）である。
\path{computations/klein_three_lines_supports.cpp} と
\path{computations/verify_klein_three_lines.py} が各段階と全代表の検査を再現する。
\end{proof}

\begin{proof}[六長方形の正規化]
残る60個の非純成分は非零である。各階数1ブロックで第零行・列を基準に
$x_{00}x_{ik}=x_{i0}x_{0k}$ ($1\le i,k\le3$) を課す。54行60列の指数行列の
Smith非零対角は $1^{46},2$、四位置の対角基底変更の階数は13である。
従って連結核は基底変更で尽くされ、二つの成分が残る。
一つの代表は全係数1であり、もう一つは
$\prod_{r<s}(-1)^{(i_r)_{\rm hi}(i_s)_{\rm lo}}$ である。
検算では左右の整可逆変換も求め、後者のSmith成分が実際に $-1$ であることを確認する。
捻りは総次数零の四語で巡回不変であり、左右のKoszul複体と双対完全対合を保存する
（\path{sections/klein_second_pairings.tex}）。従って以下は捻りを戻して調べればよい。

対称代表の階数3ブロックは両切断とも $J+\operatorname{diag}(d_a)$、$J$ は全成分1である。
$d_a$ が全て非零なら階数3の条件は $1+\sum 1/d_a=0$ で、核は $p_a=-1/d_a$ で張られる。
零の $d_a$ があるとき核はその場所に支えられ成分和が零のベクトル全体で、階数3にはちょうど二つの零が必要である。
その場合には他の位置の純平方が連続した関係になり、標準型にならない。

全 $p_a\ne0$ の場合、縮約像 $R\subset\operatorname{Sym}^2 V$ は
平方成分で $p$ に直交する超平面と、三つの色ごとの混合対称成分の和の直線からなる。
従って $v^2\in R$ には三通りの $v_0v_i=v_jv_k$ と $\sum p_i v_i^2=0$ が必要十分である。
一成分が零なら前者は座標点のみを許し、全成分が非零なら全 $v_i^2$ は等しい。
$p_i\ne0$ と $\sum p_i=1$ により、どちらにも射影解はない。
従って $R$ の対称二次直交補空間は基点自由な四式を与える。
定理\ref{thm:symmetric-frobenius-chart}の二次完全交叉であり、そのsocleテンソルは
四式で消される非零四次形式 $W$ と比例する。四位置の巡回不変性は捻った後にも成り立つ。
\end{proof}

\begin{proof}[反対階数の特殊化]
対称形では全切断が同階数である。捻った形の反対切断は
$(i,i+b)\mid(j,j+b)$ の各非零色で行列式の絶対値16の符号行列を与える。
色零では $H-\operatorname{diag}(1/p_i)$ となる。$H$ は対称で
$H^2=4I$、対角成分は1、二次主小行列式は0、三次主小行列式は全て $-4$、行列式は $-16$ である。
$P=\operatorname{diag}(p_i)$ と置くと
\[
 \det(I-PH)=1-\sum p_i+4\sum_{i<j<k}p_i p_j p_k-16\prod p_i
 =4\prod p_i\left(\sum_i1/p_i-4\right).
\]
階数2以下なら第 $l$ 行列を除く三次主小行列式も零であり
$p_l+4\prod p_i/p_l=0$ となる。従って全 $p_l$ は一つの $a$ の符号倍である。
負号の個数を $m$ とすれば $(4-2m)a=1$ かつ $1=-4(-1)^m a^2$ なので、
可能なのは上記四点だけである。各点の色零ブロックに対して直接階数2を得る。
従って反対階数は記載の三値で尽くされ、二表示の階数集合は交わらない。
\end{proof}
